\documentclass[a4paper]{article}
\usepackage[pdftex]{graphicx}
\usepackage[utf8]{inputenc}
\usepackage{enumerate}
\usepackage{icomma}
\usepackage{amssymb}
\usepackage{tikz}
\usepackage{href-ul}
\hypersetup{
	colorlinks=true,
	linkcolor=blue,
	urlcolor=blue}
\usepackage{geometry}
\geometry{a4paper, top=15mm, left=15mm, right=15mm, bottom=15mm,
	headsep=10mm, footskip=12mm}

\begin{document}
	\begin{tikzpicture}(10,3)
		\draw[ultra thick](2,0) --(10,0) -- (10,1.3) --(2,1.3) -- (2,0);
		\draw[fill=black](2,0)-- (0,.6) -- (2,.6) -- (2,0);
		\node at (6,.6) {\large Solution sheet from \href{http://www.okuyakl.com}{www.okuyakl.com}};
	\end{tikzpicture}
	\vspace{0.5 cm}
	
	\noindent{\bf Task 1. a)}\\
	We can exclude and get as double zero:
	$$x^2 ({3\over4 }x^2-4x +6)=0 \quad \Rightarrow \quad x_{1/2}=0$$
	Examination of the quadratic term in the brackets yields for the discriminant $D$:
	$$D=b^2-4ac=4^2-4\cdot {3\over 4}\cdot 6 = -2 < 0$$
	There are no other zeros.
	
	\noindent{\bf Task 1. b)}\\
	We derive and set equal to zero:
	$$
	\renewcommand{\arraystretch}{2}
	\begin{array}{rcll}
		f'(x)=0&=&3x^3-12x^2+12x \\
		0&=&3x(x^2-4x+4) & x_1=0 \\
		0&=&(x-2)^2 & x_{2/3}=2
	\end{array}
	$$
	We get the coordinates of the found locations by substituting them into $f(x)$:
	$$f(0)=0 \quad \Rightarrow \quad P_1(0|0)$$
	$$f(2)={3\over 4}\cdot 16 - 4\cdot 8 +6\cdot 4 = 4 \quad \Rightarrow \quad P_2(2|4)$$
	
	\noindent{\bf Task 2.}\\
	General functional term of the 3rd degree:
	$$f(x)=ax^3+bx^2+cx+d$$
	A fully rational function is point-symmetric if it only contains odd powers of $X$:
	$$\Rightarrow \quad f(x)=ax^3+cx$$
	It contains the point $H(-1|4)$:
	$$f(-1)=-a-c=4 \quad I $$
	At the high point the derivative is zero:
	$$f'(x)=3ax^2 + c \quad \Rightarrow \quad 3a +c = 0 \quad II $$
	We substitute I into II and get:
	$$ 3a -4 -a =0 \quad \Rightarrow \quad a = 2 \quad \Rightarrow \quad c = -6 $$
	The function term is:
	$$f(x)= 2x^3 -6x$$
	
	\noindent{\bf Task 3. a)}\\
	$$f'(x)=3x^5-2.5 + 4 \cdot x^{-3}$$
	
	\noindent{\bf Task 3. b)}\\
	$$f'(x)=(2x^{2.5})' = 5x^{1.5}$$
	
	\noindent{\bf Task 4. a)}\\
	We derive:
	$$f'(x)= 0.75 x^2 + 6x +9 $$
	We get the tangent slope at point $B(-4|-4)$ by substituting:
	$$ m = f'(-4)= 12 - 24 +9 = -3$$
	The angle you are looking for satisfies the relationship:
	$$ m = \tan\alpha \quad \Rightarrow \quad \alpha =\tan^{-1}(-3) = -71^\circ = 109^\circ$$
	
	\noindent{\bf Task 4. b)}\\
	The function has two tangents to this slope if there are two solutions to the equation:
	$$
	\renewcommand{\arraystretch}{2}
	\begin{array}{rcll}
		0.75 x^2 + 6x +9 &=& -3 \\
		0.75 x^2 +6x +12 &=& 0 \\
		x_{1/2} &=& {-6 \pm \sqrt{36-4\cdot 0.75 \cdot 12}\over 1.5}\\
		x_1 = x_2 &=& -{6 \over 1.5}= - 4
	\end{array}
	$$
	There is only one solution; so there is no other parallel tangent. (For those interested: the tangent goes through the turning point)
	
	\noindent{\bf Task 5.}\\
	First we derive and insert the two zeros:
	$$f'(x)=-2x + 5 \quad \Rightarrow \quad f'(0)= 5 \qquad f'(5)=-5$$
	
	The associated normal slopes are the negative reciprocals:
	$$m_1 = -{1 \over 5} \qquad m_2 = {1\over 5}$$
	
	This means that the two function terms of the normals are:
	$$ 0 = t_1 \quad \Rightarrow \quad n_1: y= -{1 \over 5}x $$
	$$ 0 = {1\over 5} \cdot 5 + t_2 \quad \Rightarrow \quad t_2 = -1 \quad \Rightarrow \quad n_2: y={1\over 5}x-1 $$
	
	We equate both lines and get the intersection point:
	$$ -{1 \over 5}x = {1\over 5}x-1 \quad \Rightarrow \quad x_S = 2.5 \quad \Rightarrow \quad y_S = -0.5 \quad \Rightarrow \quad P (2.5|-0.5)$$
	
	The triangle you are looking for is isosceles with the base on the x-axis, the base length is 5. The height of the triangle is the y-coordinate of P, i.e. -0.5. The triangle area is the amount of:
	$$ A = {1\over 2} \cdot 5 \cdot 0.5 = 1.25~FE$$
	
	\noindent{\bf Task 6. a)}\\
	A possible function term is:
	$$f(x)=-x^2+9 $$
	
	
	
	\begin{minipage}{0.5\textwidth}
		\noindent{\bf Task 6. b)}\\
		\includegraphics[width= 7 cm]{pavier060}
	\end{minipage}
	\begin{minipage}{0.45\textwidth}
		\noindent{\bf Task 6. c)}\\
		The rectangle has width $2u$ and height:
		$$f(u)=-u^2+9$$
		The area is the product of both, so:
		$$A(u)=2u\cdot(-u^2+9) = -2u^3 + 18u $$
		To get the maximum, we derive and set equal to zero:
		$$A'(u)=-6u^2+18 =0 \Rightarrow u_{1/2}=\pm \sqrt{3}$$
		For $u=\sqrt{3}$ the area becomes maximum.
	\end{minipage}
	
	\begin{center}
		\includegraphics[width=7 cm]{../../viecher/eendcomic.pdf}
		
		Here you can go back to the \href{https://www.okuyakl.de/math/m11garaL060/ae060.pdf}{task sheet}
	\end{center}
	
	
\end{document}